% =============================================================================
% Kapitel 4: Methoden
% -----------------------------------------------------------------------------
% Was hinein gehört (Richtwert 12--18 Seiten):
%  4.1 Problemformulierung
%      - Eingabe: Hologramm I (ggf. flat-field-korrigiert, normiert, Größe NxN)
%      - Ausgabe: z01 (oder Fr, oder log Fr) -- Begründung der Wahl, Umrechnung
%        über Gl. (eq:grundlagen:fresnelzahl); Metriken: MAE z01 [mm], rel. Fr-Fehler [%]
%  4.2 Datenerzeugung mit HoloWizard/HoloForge
%      - Simulationspipeline: Phantome (Art, Material delta/beta, Größen),
%        Beleuchtung/Probe, Geometrie (z01-Bereich, z02, Pixelgröße, Energie),
%        Rauschen (Photonenstatistik, Detektor), Flat-Field
%      - Datensatz-Versionen: configs/data_small.yaml (CPU/Smoke), configs/data_p05.yaml (GPU)
%      - Splits train/val/test, Speicherformat HDF5 + meta.json (Hash notieren!)
%      - Tabelle der Simulationsparameter (Beispiel unten)
%  4.3 Modelle
%      - Baseline: CNN-Regressor (Architektur-Skizze, Parameterzahl)
%      - Varianten: ResNet-Backbone, Fourier-Eingabe (|FFT|), Multi-Skalen,
%        Klassifikation + Verfeinerung
%      - SBI: NPE mit Normalizing Flow (sbi-Toolkit), Prior über Fr
%  4.4 Training
%      - Verlust, Optimierer, LR-Schedule, Batchgröße, Epochen, Early Stopping,
%        Augmentierung, Seeds; Hardware (CPU lokal / Colab T4 / Cluster)
%      - Tabelle Hyperparameter (aus configs/base.yaml generieren)
%  4.5 Evaluation
%      - Metriken, Referenz: modellbasierter Autofokus (HoloWizard) als Baseline
%      - Laufzeitmessung (Inferenz vs. Optimierung), Protokoll
%  4.6 Implementierung & Reproduzierbarkeit
%      - Repo-Struktur, CLI (src.data.generate_data / src.train / src.evaluate),
%        Experiment-Logbuch, Commit-SHA je Ergebnis, Software-Zitate
% =============================================================================
\chapter{\lang{Methoden}{Methods}}
\label{ch:methoden}

\section{\lang{Problemformulierung}{Problem Formulation}}
\label{sec:methoden:problem}

% Beispiel: Gleichung mit argmin und Erwartungswert
\lang{Gesucht ist eine Abbildung $f_\theta: \R^{N\times N} \to \R$, die aus einem Hologramm $I$ den Abstand $\zOne$ schätzt. Die Parameter $\theta$ werden durch Minimierung des erwarteten Verlusts bestimmt:}%
{We seek a mapping $f_\theta: \R^{N\times N} \to \R$ that estimates the distance $\zOne$ from a hologram $I$. The parameters $\theta$ are obtained by minimising the expected loss:}
\begin{equation}
  \theta^\ast = \argmin_\theta\; \E_{(I,\zOne)\sim p_{\mathrm{sim}}}\big[\,\ell\big(f_\theta(I),\, \zOne\big)\big],
  \qquad
  \ell(\hat z, z) = \abs{\hat z - z}.
  \label{eq:methoden:ziel}
\end{equation}
\todo{\lang{Wahl der Zielgröße ($\zOne$ vs.\ $\Fr$ vs.\ $\log\Fr$) begründen.}{Justify the choice of target ($\zOne$ vs.\ $\Fr$ vs.\ $\log\Fr$).}}

\section{\lang{Datenerzeugung mit HoloWizard}{Data Generation with HoloWizard}}
\label{sec:methoden:daten}

\lang{%
Die Trainingsdaten werden mit dem Simulationsmodul von HoloWizard
\cite{dora2026holowizard} erzeugt; das zugrunde liegende Vorwärtsmodell entspricht
\cref{eq:grundlagen:propagation}. \Cref{tab:methoden:simparams} fasst die
Simulationsparameter zusammen.
}{%
Training data are generated with the simulation module of HoloWizard
\cite{dora2026holowizard}; the underlying forward model corresponds to
\cref{eq:grundlagen:propagation}. \Cref{tab:methoden:simparams} summarises the
simulation parameters.
}

% Beispiel: Tabelle mit siunitx-Spalten (S) und Einheiten in der Kopfzeile
\begin{table}[htbp]
  \centering
  \caption{\lang{Simulationsparameter des Datensatzes \texttt{data\_small} (Platzhalterwerte; aus \texttt{meta.json} übernehmen).}{Simulation parameters of the \texttt{data\_small} dataset (placeholder values; copy from \texttt{meta.json}).}}
  \label{tab:methoden:simparams}
  \begin{tabular}{@{}l S[table-format=2.3] S[table-format=2.3] l@{}}
    \toprule
    \lang{Parameter}{Parameter} & {\lang{min}{min}} & {\lang{max}{max}} & \lang{Einheit}{Unit} \\
    \midrule
    $\zOne$                      & 0.020 & 0.500 & \unit{\metre} \\
    $\zTwo$                      & 20.000 & 20.000 & \unit{\metre} \\
    \lang{Photonenenergie}{Photon energy} & 11.000 & 11.000 & \unit{\kilo\electronvolt} \\
    \lang{Pixelgröße}{Pixel size} $\dx$ & 6.500 & 6.500 & \unit{\micro\metre} \\
    \lang{Bildgröße}{Image size} & {128} & {128} & \lang{Pixel}{pixels} \\
    \bottomrule
  \end{tabular}
\end{table}

\todo{\lang{Phantome, Rauschmodell, Flat-Field, Splits und \texttt{meta.json}-Hash dokumentieren.}{Document phantoms, noise model, flat field, splits and the \texttt{meta.json} hash.}}

\section{\lang{Modelle}{Models}}
\label{sec:methoden:modelle}
\todo{\lang{Baseline-\ac{cnn}, Varianten, \ac{npe}-Modell \cite{boelts2025sbireloaded}.}{Baseline \ac{cnn}, variants, \ac{npe} model \cite{boelts2025sbireloaded}.}}

\section{Training}
\label{sec:methoden:training}
\todo{\lang{Hyperparameter-Tabelle aus \texttt{configs/base.yaml}; Seeds; Hardware.}{Hyperparameter table from \texttt{configs/base.yaml}; seeds; hardware.}}

\section{Evaluation}
\label{sec:methoden:evaluation}

\lang{Als Hauptmetrik dient der \ac{mae} des Abstands in Millimetern,}{The primary metric is the \ac{mae} of the distance in millimetres,}
\begin{equation}
  \MAE_{\zOne} = \frac{1}{K}\sum_{k=1}^{K}\abs{\hat z_{01}^{(k)} - z_{01}^{(k)}},
  \label{eq:methoden:mae}
\end{equation}
\lang{ergänzt um den relativen Fehler der Fresnel-Zahl $\abs{\hat\Fr - \Fr}/\Fr$ in Prozent.}{complemented by the relative error of the Fresnel number $\abs{\hat\Fr - \Fr}/\Fr$ in per cent.}

\section{\lang{Implementierung und Reproduzierbarkeit}{Implementation and Reproducibility}}
\label{sec:methoden:implementierung}
\todo{\lang{Repo-Struktur, CLI-Befehle, Logbuch, Commit-SHAs; Software: PyTorch \cite{paszke2019pytorch}, HoloWizard \cite{dora2026holowizard}, sbi \cite{tejerocantero2020sbi}.}{Repository structure, CLI commands, logbook, commit SHAs; software: PyTorch \cite{paszke2019pytorch}, HoloWizard \cite{dora2026holowizard}, sbi \cite{tejerocantero2020sbi}.}}
